\begin{center}
{\zihao{3}\bfseries 摘\quad 要}
\end{center}

动力电池包是新能源汽车制造中装配约束强、检测要求高的关键产品，其制造过程同时涉及结构装配、电连接、热管理集成和安全测试等任务。现有动力电池包产线虽然已经引入机器人、机器视觉和数字化追溯等技术，但装配与检测往往仍以分散方式部署，工艺知识、设备能力与质量数据之间缺少统一关联，因此难以满足多品种、小批量、高一致性制造对柔性和闭环决策的要求。

本文以动力电池包这一具体产品为对象，围绕装配工艺过程、装配过程中采用的数字化手段与装备、关键智能技术及新型装配模式展开分析。首先，结合动力电池包的产品结构，梳理其装配任务、关键工艺链和主要检测节点；其次，分析机器人、人机协作、机器视觉、AI、知识图谱和大模型智能体等多类技术在装配现场中的作用边界及协同机制；最后，提出一种知识驱动的人机协同动力电池包智能装配与在线检测模式。该模式以机器人、人机协作和机器视觉构成执行层，以工艺知识图谱、规则推理和大模型智能体构成认知辅助层，强调工艺约束的显式表达、在线检测的闭环反馈以及异常诊断的辅助决策。

分析表明，动力电池包智能装配的关键不在于叠加单一前沿算法，而在于将工艺知识、设备能力、检测结果与人员决策整合为统一系统。所提出模式在换型效率、异常响应和质量追溯方面具有应用潜力，但其工程落地仍受知识维护成本、工业数据治理、模型可信性和安全认证等因素制约。总体上，本文围绕具体产品、工艺过程、系统构成、关键技术路径与应用评价，形成了较完整的论述框架。

\vspace{1em}
\noindent\textbf{关键词：}动力电池包；智能装配；在线检测；知识图谱；大模型智能体
